\title{FlashAttention-3: Fast and Accurate Attention with Asynchrony and Low-precision}

\begin{abstract}
As a foundational component within the Transformer architecture, attention mechanisms constitute the principal computational constraint in scaling large language models and supporting long-context scenarios. \textsc{FlashAttention} introduced a GPU-based acceleration strategy by curtailing memory access. Nevertheless, this approach does not exploit the advanced features available on next-generation hardware platforms; notably, \textsc{FlashAttention-2} reaches only 35\% resource utilization on the H100 GPU. To address this limitation, we present three strategies for optimizing attention on Hopper GPUs: harnessing the asynchronous capabilities of both Tensor Cores and TMA by (1) implementing warp-specialization to achieve concurrent computation and data transfer, (2) interleaving block-wise matrix multiplication with softmax calculations, and (3) employing block quantization and incoherent processing informed by hardware-enabled FP8 precision. Using these innovations, our proposed method, \textsc{FlashAttention-3}, accelerates attention computation on H100 GPUs by a factor of 1.5–2.0$\times$, delivering up to 740 TFLOPs/s (75\% efficiency) in FP16, and approaching 1.2 PFLOPs/s using FP8. Empirically, we find that FP8 \textsc{FlashAttention-3} produces a 2.6$\times$ reduction in numerical error in comparison to standard FP8 attention implementations.
\end{abstract}

\section{Introduction}
\label{sec:intro}

Within the Transformer architecture~\citep{vaswani2017attention}, the central computational limitation arises from the attention mechanism, as calculating self-attention between query and key pairs incurs a quadratic cost relative to sequence length. Overcoming this bottleneck to enable attention over extended contexts has the potential to facilitate advances in various domains, such as the analysis and reasoning across numerous lengthy documents~\citep{guo2021longt5,shaham2022scrolls,peng2023yarn} and the processing of large codebases~\citep{roziere2023code, li2023starcoder}, as well as supporting novel modalities like high-resolution visual content~\citep{chen2022scaling}, audio data~\citep{gulati2020conformer}, and video streams~\citep{ho2022video}. It can further empower new application domains, including handling user interactions characterized by lengthy historical context~\citep{sun2019bert4rec} and managing agent workflows spanning extensive time horizons~\citep{yao2022react}. Accordingly, there has been pronounced momentum towards accelerating attention for long-context settings through various means, such as efficient approximation methods~\citep{katharopoulos2020transformers,choromanski2020rethinking, tay2020efficient}, improvements in software implementation (\citep{rabe2021self,dao2022flashattention,kwon2023efficient}), and even the design of fundamentally new architectures~\citep{peng2023rwkv,sun2023retentive,gu2023mamba}.

This study extends the contributions of \citet{dao2022flashattention}, who pioneered exact-attention algorithms that strategically leverage the architectural and execution features of GPUs within their algorithmic frameworks. Dao et al. presented \textsc{FlashAttention} in \citep{dao2022flashattention}, introducing an innovative tiling method that achieves efficient attention computation by combining all operations into a single GPU kernel, thereby removing the need for intermediate accesses to global memory, which are typically a performance bottleneck.

Building on this, \citet{dao2023flashattention2} reformulated the approach as \textsc{FlashAttention-2}. This version incorporated parallelization along the sequence length and executed the primary loop of the forward pass across blocks of the key and value matrices, leading to improvements in both workload distribution and GPU occupancy. Nonetheless, our findings reveal that \textsc{FlashAttention-2} demonstrates lower resource utilization on current GPU architectures when compared to highly-tuned matrix multiplication (GEMM) kernels, achieving only about 35\% utilization compared to the 80-90\% seen on the Hopper H100 GPU. This gap may stem in part from practical implementation aspects, such as the absence of Hopper-specific instructions, which remain unchanged from those used for Ampere Tensor Cores. Recent work, including ThunkerKitten~\citep{spector2024thunder} and cuDNN 9~\citep{cudnn9}, has indicated that introducing Hopper-tailored instructions and tile-level abstractions can accelerate attention mechanisms as well as simplify code design.

At its core, the algorithm underpinning \textsc{FlashAttention-2} operates according to a straightforward synchronous paradigm, without directly leveraging asynchronous execution or reduced-precision arithmetic in its architecture. Asynchronous processing arises from specialized hardware components that target the acceleration of critical machine learning tasks: dedicated matrix multiplication units (Tensor Cores) and specialized memory access engines (Tensor Memory Accelerator -- TMA), operating independently from general-purpose CUDA cores responsible for logic, integer, and floating-point computations. The adoption of low-precision formats—such as FP8 in Hopper and FP4 in Blackwell, following the trajectory established by FP16 (introduced with Pascal in 2017) and BF16 (featured in Ampere in 2020)—demonstrates empirically that throughput can be doubled or even quadrupled while maintaining the same energy consumption and silicon footprint. A comprehensive survey of Hopper’s enhancements regarding these functionalities is provided in \cref{subsec:hardware}. 

The central engineering task is thus to modify \textsc{FlashAttention-2} so that it can exploit these hardware capabilities. Asynchronous execution calls for concurrent operation between matrix multiplication and softmax computation, despite the sequential dependency between them; meanwhile, accommodating low-precision computation entails careful handling of quantization errors, particularly to manage outlier features in large language models~\citep{dettmers2208llm, sun2024massive}.

Accordingly, we introduce \textsc{FlashAttention-3}, which integrates and advances three innovative techniques to enhance efficiency on modern GPU platforms.\footnote{While our experiments focus on NVIDIA's Hopper architecture, the proposed algorithm is compatible with any GPU system that offers strong support for asynchronous operations and low-precision computation.}

\begin{enumerate}

\item \textbf{Producer-Consumer asynchrony:} We propose a warp-specialized software pipelining strategy that leverages the decoupled execution patterns of data transfers and Tensor Cores by assigning producer and consumer tasks to distinct warps. This separation enhances the algorithm’s capability to mask both memory access and instruction issue delays.
\item \textbf{Hiding softmax under asynchronous block-wise GEMMs:} We simultaneously execute the less computationally intensive operations necessary for softmax—such as floating point multiply-add and exponentiation—while utilizing asynchronous WGMMA instructions for GEMM. To accomplish this, we modify the \textsc{FlashAttention-2} algorithm in order to eliminate certain serial dependencies between softmax and GEMM stages. In the 2-stage variant of our approach, for instance, softmax processes a current block of the score matrix concurrently with WGMMA, which calculates the next block in the background.
\item \textbf{Hardware-accelerated low-precision GEMM:} We reconfigure the forward pass to support execution on FP8 Tensor Cores for GEMM, nearly doubling the observed TFLOPs/s performance. This adaptation necessitates addressing discrepancies in memory layout expectations between WGMMA’s FP32 accumulator blocks and FP8 operand matrices. To reduce the precision-related accuracy loss inherent to FP8 computations, we implement block quantization and incoherent processing methods.

In order to assess the empirical effectiveness of our approach, we conduct a series of benchmarks of \textsc{FlashAttention-3} on the H100 SXM5 GPU, varying relevant parameters. Our results indicate that (1) the FP16 variant provides a 1.5–2.0$\times$ acceleration compared to \textsc{FlashAttention-2} in the forward computation (reaching up to 740 TFLOPs/s) and achieves a 1.5–1.75$\times$ gain in the backward computation, (2) FP8 surpasses 1.2 PFLOPs/s in throughput, and (3) when sequence length is large, FP16 consistently exceeds the performance of cuDNN's attention implementation, while FP8 also remains competitive\footnote{Specifically, with a head dimension of 64, FP8 \textsc{FlashAttention-3} is superior, whereas for head dimensions 128 and 256 it matches cuDNN attention if no causal masking is applied and trails when causal masking is present.}. Furthermore, we confirm that FP16 \textsc{FlashAttention-3} incurs the same numerical error as \textsc{FlashAttention-2}, and outperforms conventional attention mechanisms due to maintaining intermediary calculations (such as softmax scaling) in FP32 precision. Additionally, FP8 \textsc{FlashAttention-3} employing block quantization and incoherent computation demonstrates 2.6$\times$ higher numerical fidelity compared to standard attention utilizing per-tensor quantization, especially in scenarios involving outlier features.

We have released \textsc{FlashAttention-3} under a flexible license\footnote{\textsc{FlashAttention-3}
  can be found at \texttt{https://github.com/Dao-AILab/flash-attention}}, and intend to incorporate its functionalities into both PyTorch and Hugging Face libraries in order to maximize its accessibility for researchers and practitioners.

\section{Background: Multi-Head Attention and GPU Characteristics}
\label{sec:background}

\subsection{Multi-Head Attention}
\label{subsec:multi_head_attn}

Let $\mathbf{Q}, \mathbf{K}, \mathbf{V} \in \mathbb{R}^{N \times d}$ be the query, key and value input sequences associated to a single head, where $N$ is the sequence length and $d$ is the head dimension. Then the attention output $\mathbf{O}$ is computed as:

\begin{equation*}
  \mathbf{S} = \alpha \mathbf{Q} \mathbf{K}^\top \in \mathbb{R}^{N \times N}, \quad \mathbf{P} = \mathrm{softmax}(\mathbf{S}) \in \mathbb{R}^{N \times N}, \quad \mathbf{O} = \mathbf{P}\mathbf{V} \in \mathbb{R}^{N \times d},
\end{equation*}

In this formulation, $\mathrm{softmax}$ is performed along each row, and the scaling parameter $\alpha$ is often chosen as $1/\sqrt{d}$. To address potential numerical instability when computing exponentials, it is standard to subtract $\mathrm{rowmax}(\mathbf{S})$ from $\mathbf{S}$. For multi-head attention (MHA), distinct query, key, and value projection matrices are assigned to each head. This process is efficiently parallelized across the various heads and data batches, resulting in the assembled output tensor.

Now let $\phi$ be a scalar loss function and let $\mathbf{d}(-) = \partial \phi / \partial (-)$ be notation for the gradient. Given the output gradient $\mathbf{dO} \in \mathbb{R}^{N \times d}$, we compute $\mathbf{dQ}$, $\mathbf{dK}$, and $\mathbf{dV}$ according to the chain rule as follows:

\begin{align*}
  \mathbf{dV} &= \mathbf{P}^\top \mathbf{dO} \in \mathbb{R}^{N \times d} \\
  \mathbf{dP} &= \mathbf{dO} \mathbf{V}^\top \in \mathbb{R}^{N \times N} \\
  \mathbf{dS} &= \mathrm{dsoftmax} (\mathbf{dP}) \in \mathbb{R}^{N \times N} \\
  \mathbf{dQ} &= \alpha \mathbf{dS} \mathbf{K} \in \mathbb{R}^{N \times d} \\
  \mathbf{dK} &= \alpha \mathbf{dS}^\top \mathbf{Q} \in \mathbb{R}^{N \times d},
\end{align*}

In this context, $\mathbf{d}s = (\mathrm{diag}(p) - p p^\top)\mathbf{d}p$ expresses the relationship for $p = \mathrm{softmax}(s)$, where $s$ is a vector. The notation $\mathrm{dsoftmax}(\mathbf{dP})$ denotes the application of this expression to each row individually. For the backward propagation step in MHA, these calculations can once more be carried out concurrently over all heads and batch elements.

\subsection{GPU hardware characteristics and execution model}
\label{subsec:hardware}

We outline components of the GPU execution model pertinent to \textsc{FlashAttention-3}, specifically examining the NVIDIA Hopper architecture as a representative example.

\paragraph{Memory hierarchy:} GPU memory structures are designed as layered data storage, where bandwidth increases as capacity decreases (\cref{tab:gpu-hierarchy})\footnote{\citet{luo2024benchmarking} estimates shared memory throughput at 128 bytes per clock cycle per SM; this value is aggregated across 132 SMs and scaled by the boost clock rate of 1830 MHz.}. The main external memory, termed global memory (GMEM) or HBM, is an off-chip DRAM resource accessible by all streaming multiprocessors (SMs). Information retrieved from GMEM is automatically cached in a dedicated on-chip L2 cache. Each SM is equipped with a modest-sized, programmer-controlled and banked shared memory (SMEM) on-chip. At the deepest level, each SM features a register file.

\paragraph{Thread hierarchy:} Within the GPU programming framework, execution units are logically grouped into structures referred to as threads. The thread hierarchy spans multiple granularities, beginning with individual threads, followed by warps (consisting of 32 threads each), warpgroups (formed by 4 sequential warps), threadblocks—also known as cooperative thread arrays (CTAs)—then threadblock clusters (specific to Hopper architectures), and culminating with grids.

The two hierarchies exhibit a strong interconnection. Threads that belong to a single CTA are executed together on a single SM, while CTAs within one cluster are grouped for execution on the same GPC. All threads inside a CTA can directly access SMEM, in contrast to RMEM, where each thread has exclusive access to up to 256 registers.

\paragraph{Asynchrony and warp-specialization:}

GPUs are designed as throughput-oriented processors, employing concurrency and asynchronous operations to mitigate latencies in memory access and computation. The Hopper architecture introduces the Tensor Memory Accelerator (TMA), a specialized hardware component for handling asynchronous memory transfers between GMEM and SMEM \cite[\S7.29]{cuda}. Additionally, in contrast to earlier designs like Ampere, Hopper's Tensor Core—accessible through the warpgroup-wide WGMMA instruction \cite[\S9.7.14]{ptx}—operates asynchronously and is capable of fetching data directly from shared memory.

The presence of hardware mechanisms for asynchrony enables kernels optimized for warp specialization, assigning warps within a CTA exclusively to either producer or consumer functions; as a result, individual warps are dedicated solely to either data transfer operations or computational tasks. This architectural feature generally enhances the compiler's capacity to produce highly efficient instruction schedules \citep{warp-specialization-2011}. Moreover, Hopper introduces the capability for warpgroups to dynamically adjust register allocations using the \verb|setmaxnreg| instruction \cite[\S9.7.17.1]{ptx}, thereby permitting warps engaged in MMAs to access a greater portion of RMEM resources compared to those warps managing TMA operations, which require only a single thread.

\paragraph{Low-precision number formats:}
\label{sec:low-precision-gpu}
Contemporary GPUs are equipped with dedicated hardware components designed to optimize low-precision calculations. As an illustration, the WGMMA operation can be utilized to engage the FP8 Tensor Cores on Hopper, enabling a throughput per SM that is twice as fast as that achieved using FP16 or BF16.

Nonetheless, proper use of FP8 WGMMA requires familiarity with the operand layout restrictions. Suppose a GEMM operation is performed as $A \times B^{\top}$, with $A$ as an $M\times K$ matrix and $B$ as an $N\times K$ matrix. The operand $A$ or $B$ is characterized as \emph{mn-major} if memory is contiguous along the outer $M$ or $N$ dimension, and as \emph{k-major} when contiguity is along the inner $K$ dimension. FP16 WGMMA accommodates both mn-major and k-major operand layouts in SMEM, whereas FP8 WGMMA imposes a stricter requirement, permitting solely the k-major arrangement. Additionally, operations such as attention—which benefit from merging consecutive GEMMs within a single kernel—encounter difficulties due to incompatible layouts between FP32 accumulators and FP8 operands, hindering the chaining of dependent FP8 WGMMAs.

Within the domain of attention mechanisms, such layout constraints necessitate specific changes to the construction of an FP8 algorithm. The detailed approach to these adjustments is outlined in \cref{sec:algofp8}.

\subsection{Standard Attention and Flash Attention}

As described in \citet{dao2022flashattention}, we refer to \textbf{standard attention} as a GPU-based approach in which the intermediate matrices $\mathbf{S}$ and $\mathbf{P}$ are explicitly stored in HBM. The \textsc{FlashAttention} method, by contrast, was developed to minimize the overhead of intermediate memory accesses by employing a localized version of softmax reduction. This design allows the attention mechanism to be executed within a unified kernel. The local softmax operations are implemented in lines \ref{code-ws:softmax_start}-\ref{code-ws:softmax_end} within the consumer mainloop of \cref{alg:flash3_wgmma_ws_only}, in conjunction with the rescaling steps applied to blocks of $\mathbf{O}$. A straightforward derivation showing that this approach correctly computes $\mathbf{O}$ can be found in \cite[\S 2.3.1]{dao2023flashattention2}.

\section{FlashAttention-3: Algorithm}
\label{sec:algo}

This section introduces the \textsc{FlashAttention-3} algorithm. We concentrate on the forward pass for clarity; details about the backward pass are provided in~\cref{sec:algo_ws_bwd}. Initially, we outline the integration of warp-specialization, accompanied by a circular SMEM buffer, into the foundation of the \textsc{FlashAttention-2} algorithm. Subsequently, we discuss leveraging the asynchronous capabilities of WGMMA to construct a two-stage GEMM-softmax pipeline with overlapped execution. Lastly, we present the necessary adjustments for FP8, which encompass both layout alignment and maintaining accuracy through block quantization and incoherent processing.

\subsection{Producer-Consumer asynchrony through warp-specialization and pingpong scheduling}
\label{sec:algo_ws}

\paragraph{Warp-specialization}
Similar to the approach taken in \textsc{FlashAttention-2}, the forward computation in \textsc{FlashAttention-3} allows for straightforward parallelization across the dimensions of batch size, attention heads, and query sequence length. Consequently, it is sufficient to present the algorithm from the perspective of a CTA, which processes a tile $\mathbf{Q}_i$ from the query matrix in order to generate the associated output tile $\mathbf{O}_i$. For clarity, the initial explanation will focus on a warp-specialization strategy that utilizes a circular SMEM buffer, omitting for now the additional complexity introduced by overlapping GEMM and softmax operations. With $d$ denoting the head dimension and $N$ representing the sequence length, the query matrix $\mathbf{Q}$ is partitioned into blocks of size $B_r$, resulting in $T_r = \lceil \frac{N}{B_r} \rceil$ segments $\mathbf{Q}_1, .., \mathbf{Q}_{T_r}$.

\begin{algorithm}[H]
    \caption{\small\label{alg:flash3_wgmma_ws_only}\textsc{FlashAttention-3} forward pass \textbf{excluding} intra-consumer overlap -- CTA perspective}
    \begin{algorithmic}[1]
\REQUIRE Matrices $\mathbf{Q}_i \in \mathbb{R}^{B_r \times d}$, $\mathbf{K}, \mathbf{V} \in \mathbb{R}^{N \times d}$ stored in HBM; key block dimension $B_c$ with $T_c = \lceil \frac{N}{B_c} \rceil$.
\STATE Create pipeline controller to orchestrate barrier synchronization using $s$-stage circular SMEM storage.
\IF {assigned to producer warpgroup}
\STATE Release a designated quantity of registers.
\STATE Perform memory fetch of $\mathbf{Q}_i$ from HBM into shared memory.
\STATE After loading, signal consumers that $\mathbf{Q}_i$ is available.
\FOR{$0 \le j < T_c$}
    \STATE Await consumption of the $(j\,\%\,s)$th SMEM stage.
    \STATE Trigger memory loads for $\mathbf{K}_j, \mathbf{V}_j$ from HBM into shared memory at buffer stage $(j\,\%\,s)$.
    \STATE After loading, notify consumers that $\mathbf{K}_j, \mathbf{V}_j$ are ready.
\ENDFOR
\ELSE \STATE Restore a specified set of registers based on consumer warp count.
\STATE Locally initialize $\mathbf{O}_i = (0) \in \mathbb{R}^{B_r \times d}$, as well as $\ell_i, m_i = (0), (-\infty) \in \mathbb{R}^{B_r}$.
\STATE Wait until $\mathbf{Q}_i$ is present in shared memory.
\FOR{$0 \le j < T_c$}
\STATE Wait for $\mathbf{K}_j$ in shared memory.
\STATE Calculate $\mathbf{S}_i^{(j)} = \mathbf{Q}_i \mathbf{K}_j^T$ (SS-GEMM), then commit and synchronize. \label{alg:ws_only_gemm1}
\STATE Record $m_i^{\mathrm{old}} = m_i$ and update $m_i = \mathrm{max}(m_i^{\mathrm{old}}, \mathrm{rowmax}(\mathbf{S}_i^{(j)}))$. \label{code-ws:softmax_start}
\STATE Compute $\widetilde{\mathbf{P}}_i^{(j)} = \mathrm{exp}(\mathbf{S}_i^{(j)} - m_i)$ and update $\ell_i = \mathrm{exp}(m_i^{\mathrm{old}} - m_i) \ell_i + \mathrm{rowsum}(\widetilde{\mathbf{P}}_i^{(j)})$. \label{code-ws:softmax_end}
\STATE Wait for $\mathbf{V}_j$ to be ready in shared memory.
\STATE Evaluate $\mathbf{O}_i = \mathrm{diag}(\mathrm{exp}(m_i^{\mathrm{old}} - m_i))^{-1} \mathbf{O}_i + \widetilde{\mathbf{P}}_i^{(j)} \mathbf{V}_j$ (RS-GEMM), commit, and synchronize. \label{alg:ws_only_gemm2}
\STATE Unlock the $(j\,\%\,s)$th pipeline stage for producers.
\ENDFOR
\STATE Finalize $\mathbf{O}_i = \mathrm{diag}(\ell_i)^{-1} \mathbf{O}_i$ and $L_i = m_i + \log(\ell_i)$.
\STATE Write $\mathbf{O}_i$ and $L_i$ into HBM as the $i$th block in $\mathbf{O}$ and $L$.
\ENDIF
\end{algorithmic}
\end{algorithm}

In our deployment of \cref{alg:flash3_wgmma_ws_only} on Hopper, we utilize \verb|setmaxnreg| to manage (de)allocation operations, TMA for fetching $\mathbf{Q}_i$ and $\{ \mathbf{K}_j, \mathbf{V}_j \}_{0 \leq j < T_c}$, and WGMMA to carry out the GEMMs within the main consumer loop; the SS or RS prefixes denote whether the primary operand originates from shared memory or the register file, respectively. To understand the execution logic of \cref{alg:flash3_wgmma_ws_only}, it is important to observe that initiating TMA loads proceeds asynchronously and is not impeded by the completion status of other loads. Additionally, the producer's main loop does not initiate waits during the initial $s$ iterations, as the buffer is still in the process of being populated.

\paragraph{Pingpong scheduling} The asynchronous execution of WGMMA and TMA, paired with warp-specialization, creates the possibility to simultaneously run softmax calculations in one warpgroup while another warpgroup processes GEMM operations. To highlight why this is beneficial, consider that, on current hardware accelerators, non-matmul operations achieve markedly lower throughput relative to matmul. For instance, the H100 SXM5 GPU delivers 989 TFLOPS for FP16 matmul, but only reaches 3.9 TFLOPS for specialized functions such as exponential\footnote{The CUDA programming guide outlines that each streaming multiprocessor (SM) can handle 16 special function operations per clock cycle. This total is derived by multiplying 16 with 132 SMs and the 1830 MHz clock rate, yielding 3.9 TFLOPS for these functions.}, which are required for computing softmax. When performing an attention forward pass in FP16 at a head dimension of 128, the quantity of matmul FLOPS is 512 times greater than exponential FLOPS, yet exponential operations lag significantly in throughput—specifically, by a factor of 256. As a result, exponentials may occupy approximately 50\% of the compute cycle, compared to matmuls. This imbalance is even more pronounced using FP8, as matmul throughput doubles but exponential throughput remains unchanged.

Because the exponential operation is executed by a distinct hardware component—the multi-function unit—it is preferable to overlap its computation with the Tensor Cores' matrix multiplication. To achieve this overlap, we introduce synchronization barriers (\texttt{bar.sync} instructions) that sequence the GEMM operations: warpgroup 1's GEMMs (GEMM1, corresponding to $\mathbf{P} \mathbf{V}$ for one iteration, and GEMM0, which is $\mathbf{Q} \mathbf{K}^\top$ for the following iteration) are completed prior to the GEMMs of warpgroup 2. Consequently, the softmax calculation in warpgroup 1 occurs concurrently with warpgroup 2's GEMM operations. Afterward, the two warpgroups exchange roles, so that warpgroup 2 executes softmax while warpgroup 1 processes GEMMs—a technique known as ``pingpong'' scheduling, as depicted in~\cref{fig:pingpong_scheduling}. While real-world implementation of this scheduling approach does not achieve the perfect alternation shown in the figure, we consistently observe it leads to notable speedups (for instance, raising FP16 forward performance with head size 128 and sequence length 8192 from 570 TFLOPS to 620–640 TFLOPS).

\paragraph{Attention variants} In the case of multi-query attention~\citep{shazeer2019fast} and grouped query attention~\citep{ainslie2023gqa}, our implementation adheres to the strategy outlined in \textsc{FlashAttention-2}, modifying the tensor access patterns so as to prevent redundant copies of $\mathbf{K}$ and $\mathbf{V}$ within HBM.

\subsection{Intra-warpgroup overlapping GEMMs and softmax}

It is possible to achieve partial overlap between certain instructions from the softmax computation and those from the GEMMs within a single warpgroup. In the following, we outline a specific method for facilitating such overlap.

In the attention algorithm, operations within the inner loop (main loop) have sequential dependencies that impede parallelization within a single iteration. For example, (local) softmax (lines \ref{code-ws:softmax_start} to \ref{code-ws:softmax_end}) relies on the output $\mathbf{S}_i^{(j)}$ of the first GEMM, while the second GEMM takes its result $\widetilde{\mathbf{P}}_i^{(j)}$ as an operand. Indeed, the wait statements in lines \ref{alg:ws_only_gemm1} and \ref{alg:ws_only_gemm2} of \cref{alg:flash3_wgmma_ws_only} serialize the execution of softmax and GEMMs. However, we can break these dependencies by pipelining across iterations through additional buffers in registers. Pursuing this idea, we propose the following two-stage\footnote{Note that the number of stages of the overlapping scheme is bounded by, but need not equal, the number $s$ of stages in the circular SMEM buffer.} GEMM-softmax pipelining algorithm:

\begin{algorithm}[H]
    \caption{\small\label{alg:flash3_wgmma}\textsc{FlashAttention-3} consumer warpgroup forward pass}
    \begin{algorithmic}[1]
      \REQUIRE  Matrices $\mathbf{Q}_i \in \mathbb{R}^{B_r \times d}$ and $\mathbf{K}, \mathbf{V} \in \mathbb{R}^{N \times d}$ stored in HBM, block size $B_c$ for keys with $T_c = \lceil \frac{N}{B_c} \rceil$.
      \STATE Assign a fixed number of registers dynamically, proportional to the count of consumer warps.
      \STATE Set up on-chip memory: initialize $\mathbf{O}_i = (0) \in \mathbb{R}^{B_r \times d}$; set $\ell_i = (0)$ and $m_i = (-\infty)$ for $\ell_i, m_i \in \mathbb{R}^{B_r}$.
      \STATE Synchronize for $\mathbf{Q}_i$ and $\mathbf{K}_0$ to be copied into shared memory.
      \STATE Calculate $\mathbf{S}_{\mathrm{cur}} = \mathbf{Q}_i \mathbf{K}_0^T$ via WGMMA, then commit and wait for completion.
      \STATE Free the buffer’s initial stage corresponding to $\mathbf{K}$.
      \STATE Using $\mathbf{S}_{\mathrm{cur}}$, update $m_{i}$, $\tilde{\mathbf{P}}_{\mathrm{cur}}$, and $\ell_{i}$; appropriately rescale $\mathbf{O}_i$.
      \FOR{$1 \le j < T_c - 1$}
        \STATE Ensure $\mathbf{K}_j$ has been transferred into shared memory.
        \label{code:mainloop_start}
        \STATE Perform matrix multiplication $\mathbf{S}_{\mathrm{next}} = \mathbf{Q}_i \mathbf{K}_{j}^T$ using WGMMA; commit the operation, but no wait required. \label{code:first_wgmma}
        \STATE Synchronize for $\mathbf{V}_{j-1}$ being available in shared memory.
        \STATE Accumulate to output: compute $\mathbf{O}_{i} = \mathbf{O}_{i} + \tilde{\mathbf{P}}_{\mathrm{cur}} \mathbf{V}_{j-1}$ utilizing WGMMA; commit operation without waiting. \label{code:second_wgmma}
        \STATE Wait for completion of WGMMA product $\mathbf{Q}_i \mathbf{K}_{j}^T$.
        \STATE Based on $\mathbf{S}_{\mathrm{next}}$, update $m_{i}$, generate $\tilde{\mathbf{P}}_{\mathrm{next}}$, and set $\ell_{i}$. \label{code:softmax}
        \STATE Wait for WGMMA computation $\tilde{\mathbf{P}}_{\mathrm{cur}} \mathbf{V}_{j-1}$ and subsequently adjust $\mathbf{O}_i$ scaling.
        \STATE Release the $(j\,\%\,s)$th and $(j-1\,\%\,s)$th stages of buffers corresponding to $\mathbf{K}$ and $\mathbf{V}$ respectively.
        \STATE Assign $\mathbf{S}_{\mathrm{cur}} \leftarrow \mathbf{S}_{\mathrm{next}}$.
        \label{code:mainloop_end}
      \ENDFOR
      \STATE Wait for $\mathbf{V}_{T_c - 1}$ to be resident in shared memory.
      \STATE Finalize output by evaluating 
      $\mathbf{O}_{i} = \mathbf{O}_{i} + \tilde{\mathbf{P}}_{\mathrm{last}} \mathbf{V}_{T_c - 1}$ through WGMMA, commit, and synchronize until finished.

\STATE Epilogue: Adjust $\mathbf{O}_{i}$ according to $m_{i}$. Determine $L_{i}$ utilizing both $m_{i}$ and $\ell_{i}$.
Store $\mathbf{O}_{i}$ and $L_{i}$ in HBM as the $i$-th segments of $\mathbf{O}$ and $L$, respectively.
\end{algorithmic}
\end{algorithm}

\cref{alg:flash3_wgmma} serves as an alternative to the consumer component described in \cref{alg:flash3_wgmma_ws_only}, thereby forming the full \textsc{FlashAttention-3} procedure for FP16 computations. WGMMA is employed as a shorthand for asynchronous GEMM. In the primary loop (from lines \ref{code:mainloop_start} through \ref{code:mainloop_end}), each second WGMMA operation within iteration $j$ (see line \ref{code:second_wgmma}) is concurrently executed with the softmax computations from the subsequent iteration $j+1$ (refer to line \ref{code:softmax}).

Although the pipeline architecture depicted earlier promises notable theoretical speedups, several real-world factors must be taken into account:

\paragraph{Compiler reordering} Although the pseudocode presents a sequential execution layout, compilers like NVCC frequently restructure instructions to achieve better optimization. Such rearrangements can disturb the intended alternation of WGMMA and non-WGMMA operations, possibly resulting in erratic outcomes or reduced benefits from pipelining. Examination of the generated SASS output reveals that NVCC does produce overlapped instruction schedules, as discussed in Section~\ref{sec:2-stage-sass}.

\paragraph{Register pressure} Achieving optimal throughput demands minimal register spilling, but implementing the 2-stage pipeline increases register requirements since intermediate values and inter-stage state must be stored. In particular, each threadblock must retain an additional $\mathbf{S}_{\mathrm{next}}$ in its register set, adding a register load of $B_r \times B_c \times \text{sizeof}(\text{float})$. This heightened register consumption competes with block size scaling—another typical optimization technique that also consumes registers. In practical deployment, these competing demands should be navigated based on empirical profiling.

\paragraph{3-stage pipelining} Building upon the earlier 2-stage method, we introduce a 3-stage strategy that aims to further interleave a second WGMMA operation with the softmax step. While this scheme can theoretically improve Tensor Core occupancy, it imposes additional register overhead to maintain another stage of intermediate data, intensifying the balance required between tile size and pipeline depth. Further technical details and quantitative assessment of the 3-stage approach are available in ~\cref{sec:3-stage}.

\subsection{Low-precision with FP8}
\label{sec:algofp8}

\textbf{Efficiency: layout transformations.} Executing the forward pass of \textsc{FlashAttention-3} using FP8 precision introduces unique difficulties related to layout compliance, which are absent when working with FP16.

To begin, it is important to observe that the input tensors $\mathbf{Q}$, $\mathbf{K}$, and $\mathbf{V}$ are usually arranged with contiguity along the head dimension. However, in order to comply with the k-major constraint imposed by FP8 WGMMA during the second GEMM, $\mathbf{V}$—specifically the portions of $\mathbf{V}$ fetched into SMEM—must be contiguous with respect to the sequence length dimension. As the TMA load is incapable of modifying which dimension is contiguous, we have two primary alternatives: (1) transpose $\mathbf{V}$ in GMEM as a preparatory procedure, or (2) apply an in-kernel transpose to the segments of $\mathbf{V}$ post-transfer to SMEM. For option (1), it is possible to either (1a) integrate the transpose operation into the epilogue of a preceding process, such as rotary embedding, or (1b) utilize an independent pre-processing transpose kernel\footnote{A well-tuned transpose kernel can deliver throughput close to the device's bandwidth~\citep{colfax_cutlass_transpose_2024}.} in order to switch the sequence length and head dimension strides. Nonetheless, method (1a) poses challenges for adoption within typical libraries, while method (1b) tends to be inefficient in scenarios limited by memory bandwidth, like inference.

For FP8 \textsc{FlashAttention-3}, we select option (2) as our strategy. During the in-kernel transpose, we utilize the LDSM (\verb|ldmatrix|) and STSM (\verb|stmatrix|) instructions. These enable a warp to cooperatively move data between SMEM and RMEM in blocks of 128 bytes. \footnote{According to the PTX specification, LDSM/STSM are primarily designed for transferring $8 \times 8$ matrices of 16-bit values \cite[\S 9.7.13.4.15-16]{ptx}, but by packing two 8-bit values together, these instructions become suitable for operations with FP8 precision. However, the transpose modes of LDSM/STSM are unable to separate packed 8-bit entries, which makes certain register shuffling between LDSM and STSM necessary to effect a tile-level transpose. The specifics of this process are omitted here.} LDSM/STSM are advantageous for their register efficiency—permitting their use directly in the producer warpgroup—and also possess the capability to adjust data layout during memory transfer by performing transpositions. Furthermore, following the initial iteration, the subsequent tile of $\mathbf{V}$ can have its transposition scheduled concurrently, overlapping with the computation of the two WGMMAs that correspond to the previous $\mathbf{V}$ tile and the current $\mathbf{K}$ tile.

Second, we note that, in contrast to FP16, the FP32 accumulator utilized by FP8 WGMMA presents a memory layout that differs from the arrangement expected for operand A when stored in registers. This is illustrated in \cref{fig:rmem_accum} and \cref{fig:rmem_operand}, which provide excerpts of these respective layouts, showing how register entries are distributed among threads sequentially. Employing byte permute instructions enables conversion of the initial WGMMA accumulator’s format into one compatible with subsequent WGMMA operations and with the layout of the $\mathbf{V}$ tile generated via in-kernel transpose. Specifically, as demonstrated in \cref{fig:rmem_accum}, the ordering is modified to
$$\{ \verb|d0 d1 d4 d5 d2 d3 d6 d7| \},$$
with this register ordering repeated every 8 bytes. Logically, this technique rearranges the columns of the $\mathbf{P}$ tile (for example, columns $0189$ are repositioned as the first four columns). To ensure WGMMA produces the intended output tile, we configure the in-kernel transpose to apply an analogous row reordering to the $\mathbf{V}$ tile.\footnote{The flexibility offered by performing an in-kernel transpose removes the need for shuffle instructions to reassign register ownership across threads, as previously discussed in~\citep{colfax_fp8_flashattention_2024}.}

\textbf{Accuracy: block quantization and incoherent processing.} The FP8 (e4m3) representation allocates only 3 bits to the mantissa while reserving 4 bits for the exponent, leading to increased numerical inaccuracies compared to FP16 or BF16. In addition, in large-scale models, the occurrence of outlier values with substantially higher magnitudes than the rest is common~\citep{dettmers2208llm, sun2024massive}, which adds further complexity to the quantization process. A prevalent strategy to address this involves applying per-tensor scaling~\citep{micikevicius2022fp8}, where each tensor (such as $\mathbf{Q}$, $\mathbf{K}$, and $\mathbf{V}$) maintains a dedicated scaling factor. In order to mitigate the accuracy loss when using FP8 for attention computation, we adopt two methodological approaches:

\begin{enumerate}

\item \textbf{Block quantization}: In this approach, we assign a single scalar per block. Specifically, for each tensor $\mathbf{Q}$, $\mathbf{K}$, and $\mathbf{V}$, the tensor is divided into blocks of size $B_r \times d$ or $B_c \times d$, each of which is quantized independently. This method can be combined with certain pre-attention operations, such as rotary embedding, without incurring additional computational overhead, since the bottleneck is the memory bandwidth of rotary embedding. Because the \textsc{FlashAttention-3} algorithm is intrinsically block-based, we are able to rescale each block of $\mathbf{S}$ to incorporate block quantization, and this does not require extra computation.

\item \textbf{Incoherent processing}: To mitigate the effects of outliers, we pre-multiply $\mathbf{Q}$ and $\mathbf{K}$ by a random orthogonal matrix $\mathbf{M}$ before FP8 quantization. The property $\mathbf{M} \mathbf{M}^\top = I$ ensures that $(\mathbf{Q} \mathbf{M}) (\mathbf{K} \mathbf{M})^\top = \mathbf{Q} \mathbf{K}^\top$, so applying $\mathbf{M}$ to both does not alter the resultant attention scores. This process effectively distributes the influence of outliers, as each component of $\mathbf{Q} \mathbf{M}$ or $\mathbf{K} \mathbf{M}$ becomes a random weighted sum of the original components, thereby minimizing quantization error. Following \citet{chee2024quip} and~\citet{tseng2024quip}, $\mathbf{M}$ is constructed as the product of random diagonal matrices whose entries are $\pm 1$ and a Hadamard matrix. This formulation allows multiplication to be executed in $O(d \log d)$ rather than $O(d^2)$ time, and it can also be seamlessly fused with the rotary embedding operation without incurring additional compute.

\end{enumerate}

Experimental results confirm that these strategies decrease numerical error by as much as 2.6$\times$, as shown in \cref{sec:numerical_error}.

\section{Empirical Validation}
\label{sec:exp}

To develop and assess the performance and precision of \textsc{FlashAttention-3}, we leverage the WGMMA and TMA abstractions provided by CUTLASS~\citep{Thakkar_CUTLASS_2023}.

\begin{itemize}

\item \textbf{Attention Benchmarking.} We assess the performance of \textsc{FlashAttention-3} by evaluating its runtime over various sequence lengths, and compare these results against the standard PyTorch implementation, \textsc{FlashAttention-2}, the Triton variant of \textsc{FlashAttention-2} (leveraging H100-specific instructions), and a cuDNN-based vendor version of \textsc{FlashAttention-2} optimized for H100 GPUs. Our analysis demonstrates that \textsc{FlashAttention-3} achieves up to a 2.0$\times$ speed advantage relative to \textsc{FlashAttention-2}, and a 1.5$\times$ improvement when compared to the Triton version. Furthermore, \textsc{FlashAttention-3} attains a peak throughput of 740 TFLOPs/s, which corresponds to 75\% of the H100 GPU’s theoretical peak TFLOPs/s.
\item \textbf{Ablation Analysis.} We verify through controlled experiments that our algorithmic refinements—specifically, warp-specialization and the GEMM-softmax pipelining—are key contributors to the observed acceleration in \textsc{FlashAttention-3}.
\item \textbf{FP8 Attention Precision.} Our evaluation indicates that applying block quantization together with incoherent processing significantly mitigates numerical errors in FP8 \textsc{FlashAttention-3}, reducing them by a factor of 2.6$\times$.
\end{itemize}

\subsection{Benchmarking Attention}
\label{subsec:benchmark_attn}

We evaluate the performance of various attention mechanisms on an H100 80GB SXM5 GPU across several configurations, specifically considering settings with and without a causal mask, and head dimensions of either 64 or 128, using FP16 inputs. The empirical results, presented in~\cref{fig:benchmark_attn_fwd} and~\cref{fig:benchmark_attn_bwd}, reveal that \textsc{FlashAttention-3} achieves approximately 1.5-2.0$\times$ acceleration over \textsc{FlashAttention-2} during the forward pass, and exhibits a 1.5-1.75$\times$ speedup in the backward pass. When benchmarked against a conventional attention implementation, \textsc{FlashAttention-3} demonstrates a dramatic improvement, offering speedups of up to 3-16$\times$. Notably, for sequence lengths in the medium to large range (at least 1k tokens), \textsc{FlashAttention-3} even outperforms the cuDNN library, a proprietary and highly optimized package designed for H100 GPUs.

\paragraph{Benchmark settings:} We vary the sequence length as 512, 1k, ..., 16k, and set batch size so that the total number of tokens is 16k. We set the hidden dimension to 2048, and head dimension to be either 64, 128, or 256 (i.e., 32 heads, 16 heads, or 8 heads). To calculate the FLOPs of the forward pass, we use:

\begin{equation*}
  4 \cdot \text{seqlen}^2 \cdot \text{head dimension} \cdot \text{number of heads}.
\end{equation*}

Causal masking reduces the computed entries by half, so we divide the original quantity by 2. For the backward pass, we estimate its FLOPs by multiplying the forward pass FLOPs by 2.5, reflecting the presence of 2 matrix multiplications during the forward operation and 5 during the backward process, factoring in the recomputation step.

In addition, we evaluate the forward pass runtime utilizing FP8 within comparable conditions. The outcomes corresponding to headdim 256 are presented in ~\cref{fig:benchmark_attn_fp8}, while the comprehensive set of results is detailed in \cref{sec:benchmark_attn_fp8_full}.

\subsection{Ablation Study: 2-Stage Pipelining Experiments}
\label{sec:2-stage-pipelining-experiments}

We conduct ablation experiments on both the 2-stage WGMMA-softmax pipelining and warp-specialization in the context of non-causal FP16 \textsc{FlashAttention-3}, employing fixed parameter values $\{\text{batch}, \text{seqlen}, \text{nheads}, \text{hdim}\} = \{ 4, 8448, 16, 128\}$. As shown in~\cref{table:ablation_pipelining}, our results demonstrate that incorporating asynchrony via warp-specialization and enabling overlap between GEMM and softmax operations yields notable performance gains, increasing throughput from 570 to 661 TFLOPs.

\subsection{Numerical Error Validation}
\label{sec:numerical_error}

Given the attention drawn to the numerical errors in \textsc{FlashAttention}~\citep{golden2024flash}, we conduct a comparison of \textsc{FlashAttention-2}, \textsc{FlashAttention-3}, and a conventional attention implementation, using an FP64 reference for evaluation. In order to replicate the presence of outlier features and activations typically observed in LLMs~\citep{dettmers2208llm, sun2024massive}, the values for $\mathbf{Q}, \mathbf{K}, \mathbf{V}$ are sampled according to the following probability distribution:

\begin{equation*}
  \mathcal{N}(0, 1) + \mathcal{N}(0, 100) \cdot \mathrm{Bernoulli}(0.001).
\end{equation*}

Specifically, each element follows a normal distribution with mean zero and unit standard deviation; however, for 0.1\% of the elements, we add a separate random variable, independently sampled from a normal distribution with standard deviation 10. The root mean squared error (RMSE) resulting from these modifications is presented in~\cref{table:numerical_error}. When using FP16, both \textsc{FlashAttention-2} and \textsc{FlashAttention-3} demonstrate 1.7$\times$ reduction in RMSE relative to the typical implementation due to storing intermediate values (such as softmax outputs) in FP32. For the FP8 baseline, per-tensor scaling is applied, the matmul accumulator operates in FP32, and the intermediate softmax values are kept in FP16. Owing to block quantization and non-coherent computation, \textsc{FlashAttention-3} in FP8 yields accuracy that is 2.6$\times$ greater than the baseline.

\section{Dicussion, Limitations, Conclusion}
\label{sec:discussion}

Through \textsc{FlashAttention-3}, we have illustrated that leveraging innovative programming paradigms and exploiting modern hardware capabilities—specifically asynchrony and reduced precision—can significantly enhance both the computational efficiency and precision of attention mechanisms. Our advancements yield a 1.5–2.0$\times$ speedup relative to \textsc{FlashAttention-2}, alongside a 2.6$\times$ reduction in FP8 numerical errors compared to traditional per-tensor quantization methods. Our current approach is subject to several constraints that we aim to resolve: advancing optimization for LLM inference, implementing a persistent kernel within the FP8 kernel,\footnote{For benchmarking purposes, FP16 \textsc{FlashAttention-3} employs a persistent kernel and load balancing, unlike FP8 \textsc{FlashAttention-3}, which lacks these features. Consequently, FP8 \textsc{FlashAttention-3} underperforms on shorter sequences and causal masking relative to FP8 cuDNN kernels.} and thoroughly investigating the implications of employing low-precision attention during extensive model training. While this study centers on Hopper GPUs, we anticipate that the methodology introduced will extend to other accelerator architectures. It is our expectation that improvements in attention as a foundational primitive will enable novel uses in scenarios involving longer contexts.

\end{document}